\chapter{The best quote on each side}
\label{ch:best}

\begin{question}
Given $g$, which quote distances $\delta^b$ and $\delta^a$ reach the two $\sup$s in the HJB equation?
\end{question}

\begin{whatfor}
\begin{itemize}
  \item Reduces each $\sup$ to maximising one function of one variable.
  \item Differentiates it, finds where the derivative changes sign, and solves.
  \item Shows that $A$ cancels: how often fills come does not change where to quote.
\end{itemize}
\end{whatfor}

\begin{setup}[Setup]
The jump terms of \cref{res:jump}. Inventory $q$ and time $t$ are held fixed, so $\Delta^bg$ and $\Delta^ag$ are fixed numbers.
\end{setup}

%=====================================================================
\section{One variable per side}

\begin{definition}[The bid objective]
\label{def:f}
For a fixed number $c$, let
\[
  f(\delta)=A\,e^{-k\delta}\Big(1-e^{-\gamma(\delta+c)}\Big).
\]
\end{definition}

\begin{result}[Each $\sup$ is a $\sup$ of $f$]
\label{res:reduce}
\[
  \begin{aligned}
    \sup_{\delta^b}\lambda(\delta^b)\big[v^b-v\big]&=(-v)\sup_{\delta}f(\delta)\quad\text{with }c=\Delta^bg,\\
    \sup_{\delta^a}\lambda(\delta^a)\big[v^a-v\big]&=(-v)\sup_{\delta}f(\delta)\quad\text{with }c=-\Delta^ag.
  \end{aligned}
\]
\end{result}

\begin{proof}
\begin{steps}
  \item The bid jump term is $(-v)\,f(\delta^b)$ with $c=\Delta^bg$.
  \\ \emph{Why:} \cref{res:jump} and \cref{def:f}.
  \item The ask jump term is $(-v)\,f(\delta^a)$ with $c=-\Delta^ag$, because $\delta^a-\Delta^ag=\delta^a+c$.
  \\ \emph{Why:} \cref{res:jump} and \cref{def:f}.
  \item $-v>0$ does not depend on the quote, and a positive constant comes out of a $\sup$.
  \\ \emph{Why:} the remark after \cref{res:jump}; $\sup(aF)=a\sup F$ for $a>0$.
\end{steps}
\end{proof}

%=====================================================================
\section{The first-order condition}

\begin{result}[Derivative of $f$]
\label{res:fprime}
\[
  f'(\delta)=A\,e^{-k\delta}\Big[-k+(k+\gamma)\,e^{-\gamma(\delta+c)}\Big].
\]
\end{result}

\begin{proof}
\begin{steps}
  \item By the product rule, with $\frac{d}{d\delta}Ae^{-k\delta}=-kAe^{-k\delta}$ and $\frac{d}{d\delta}\big(1-e^{-\gamma(\delta+c)}\big)=\gamma e^{-\gamma(\delta+c)}$,
  \[
    f'(\delta)=-kA\,e^{-k\delta}\Big(1-e^{-\gamma(\delta+c)}\Big)+A\,e^{-k\delta}\,\gamma\,e^{-\gamma(\delta+c)}.
  \]
  \emph{Why:} $(uw)'=u'w+uw'$; the chain rule for each exponential.
  \item Take out $Ae^{-k\delta}$ and collect the $e^{-\gamma(\delta+c)}$ terms:
  \[
    \begin{aligned}
      f'(\delta)&=A\,e^{-k\delta}\Big[-k+k\,e^{-\gamma(\delta+c)}+\gamma\,e^{-\gamma(\delta+c)}\Big]\\
      &=A\,e^{-k\delta}\Big[-k+(k+\gamma)\,e^{-\gamma(\delta+c)}\Big].
    \end{aligned}
  \]
  \emph{Why:} $-k\big(1-e^{-\gamma(\delta+c)}\big)=-k+ke^{-\gamma(\delta+c)}$.
\end{steps}
\end{proof}

\begin{result}[The best quote distance]
\label{res:foc}
$f$ is largest at
\[
  \delta^*=\frac1\gamma\ln\!\Big(1+\frac\gamma k\Big)-c .
\]
So the best bid and ask distances are
\[
  \delta^{b*}=\frac1\gamma\ln\!\Big(1+\frac\gamma k\Big)-\Delta^bg,\qquad
  \delta^{a*}=\frac1\gamma\ln\!\Big(1+\frac\gamma k\Big)+\Delta^ag.
\]
\end{result}

\begin{proof}
\begin{steps}
  \item $Ae^{-k\delta}>0$, so $f'(\delta)$ has the sign of the bracket $B(\delta)=-k+(k+\gamma)e^{-\gamma(\delta+c)}$.
  \\ \emph{Why:} \cref{res:fprime}; $A>0$ and an exponential is positive.
  \item $B$ is strictly decreasing in $\delta$, falls below $0$ for large $\delta$, and is above $0$ for $\delta$ far enough to the left.
  \\ \emph{Why:} $e^{-\gamma(\delta+c)}$ strictly decreases from $+\infty$ to $0$ as $\delta$ grows, and $k+\gamma>0$.
  \item So $B$, and with it $f'$, is positive to the left of one point $\delta^*$ and negative to its right. At $\delta^*$,
  \[
    -k+(k+\gamma)e^{-\gamma(\delta^*+c)}=0,\qquad\text{that is}\qquad e^{-\gamma(\delta^*+c)}=\frac{k}{k+\gamma}.
  \]
  By F5, $f$ is largest at $\delta^*$.
  \\ \emph{Why:} Steps 1 and 2; F5.
  \item Take logarithms and use $\ln\frac{k}{k+\gamma}=-\ln\frac{k+\gamma}{k}=-\ln\!\big(1+\frac\gamma k\big)$:
  \[
    \begin{aligned}
      -\gamma(\delta^*+c)&=-\ln\!\Big(1+\frac\gamma k\Big)\\
      \delta^*&=\frac1\gamma\ln\!\Big(1+\frac\gamma k\Big)-c.
    \end{aligned}
  \]
  \emph{Why:} $\ln$ is the inverse of $\exp$; divide by $-\gamma$.
  \item For the bid, put $c=\Delta^bg$; for the ask, put $c=-\Delta^ag$.
  \\ \emph{Why:} \cref{res:reduce}.
\end{steps}
\end{proof}

The equation $f'(\delta^*)=0$ is the \dfn{first-order condition}. The quotes may not cross the mid-price ($\delta\ge0$). If the formula gives a negative distance, the best allowed choice is $\delta=0$, because there $f$ is still rising.

\begin{intuition}
$A$ has disappeared from $\delta^*$. $A$ multiplies the whole of $f$, so it changes how big the best value is, but not where it is. How often fills come does not matter; what matters is how fast the fill rate falls with distance ($k$) against how much the dealer dislikes risk ($\gamma$).
\end{intuition}

\begin{example}
For RELIANCE, $\gamma=0.005$ per rupee and $k=2.52$ per rupee, so the common part of both distances is
\[
  \frac1\gamma\ln\!\Big(1+\frac\gamma k\Big)=200\,\ln(1.001984)\approx0.396\text{ rupees}.
\]
A dealer with no risk correction ($\Delta g=0$) would quote about $0.40$ rupees on each side of the mid-price.
\end{example}

\begin{recall}
\textbf{Objective.} $f(\delta)=Ae^{-k\delta}\big(1-e^{-\gamma(\delta+c)}\big)$, with $c=\Delta^bg$ (bid) or $c=-\Delta^ag$ (ask).

\textbf{Derivative.} $f'=Ae^{-k\delta}\big[-k+(k+\gamma)e^{-\gamma(\delta+c)}\big]$, positive and then negative.

\textbf{First-order condition.} $e^{-\gamma(\delta^*+c)}=\frac{k}{k+\gamma}$, so
\[
  \delta^{b*}=\edge-\Delta^bg,\qquad \delta^{a*}=\edge+\Delta^ag .
\]
$A$ cancels.
\end{recall}
